\documentclass[10pt,a4paper]{amsart}
\usepackage[margin=19mm]{geometry}
\usepackage{amsmath,amssymb,mathtools}
\usepackage[scr=rsfs,cal=euler]{mathalfa}
\usepackage{xfrac}
\usepackage[unicode,hidelinks,bookmarks=false]{hyperref}
\newcommand{\ie}{i.e.\ }
\newcommand{\cf}{cf.\ }
\newcommand{\resp}{respectively\ }
\newcommand{\Subsection}{\subsection}
\providecommand{\nobreakdash}{\nobreak\hbox{-}}
\setlength{\emergencystretch}{2em}
\multlinegap0pt
\usepackage{bm}
\usepackage{bbm}
\usepackage{dsfont}
\usepackage{xspace}
\usepackage{cancel}
\usepackage{mathtools}
\usepackage{amsthm}
\makeatletter
\@ifundefined{theorem}{\newtheorem{theorem}{Theorem}}{}
\@ifundefined{claim}{\newtheorem{claim}{Claim}}{}
\@ifundefined{proposition}{\newtheorem{proposition}[theorem]{Proposition}}{}
\makeatother
\newcommand{\cO}{\mathcal{O}}
\newcommand{\cS}{\mathcal{S}}
\newcommand{\cV}{\mathcal{V}}
\title[Reconstructing a shellable sphere from its facet-ridge graph]{Reconstructing a shellable sphere from its facet-ridge graph}
\author{Yirong Yang}
\date{Source: arXiv:2401.04220v2 (CC BY 4.0)}
\begin{document}
\maketitle

\noindent\emph{Source and licence.} Reconstructing a shellable sphere from its facet-ridge graph. Version arXiv:2401.04220v2 (CC BY 4.0), distributed under the Creative Commons Attribution 4.0 International licence, as recorded in the arXiv licence metadata for this exact version and shown on the abstract page https://arxiv.org/abs/2401.04220v2. Full rights notice: licenses/arxiv-2401.04220-CC-BY-4.0.txt.\par\smallskip

\noindent\textit{Editorial scope.} Counted unit: the Claim whose source label is \texttt{clm:initial} in the article (source0.tex lines 420--425, source label \texttt{clm:initial}), delivered together with the article's own complete original proof of it (source0.tex lines 427--447, one \texttt{proof} environment of 21 lines). No mathematics is written, repaired or re-derived: statement and proof are the article's own text line for line. Required context retained uncounted: prop:reconstructshelling (lines 222--224). Every definition, display and label of those lines is retained unchanged. Omitted from this package and not redistributed: 1--221, 225--419, 426--426, 448--680. Nothing inside a retained excerpt is omitted or altered: every retained body is delivered exactly as the article's lines are written, so no omission receipt of any kind accompanies this package. Citations: the retained excerpts carry no citation command at all, so no bibliography entry of the article is transcribed and no reference is redistributed. Reference adaptation: none was needed and none was made; every reference inside the delivered unit resolves inside it, so no unresolved reference or citation is produced. Printed slips kept literal: no printed word, symbol, hypothesis or number of the counted statement or of its proof was altered. Layout and class: the article's own class is \texttt{article} with packages including inputenc, amsmath, amsthm, amssymb, geometry, enumerate, enumitem, mathrsfs, mathtools, tikz-cd, xcolor, graphicx, array, bm; the wrapper is the lane's pinned \texttt{amsart} editorial wrapper at 10pt on A4, which restates the theorem-like environments the retained text needs and adds the article's own macro definitions that the retained text needs (\texttt{\textbackslash cO}, \texttt{\textbackslash cS}, \texttt{\textbackslash cV}), with extra wrapper packages (bm, bbm, dsfont, xspace, cancel, mathtools). The delivered numbering is the unit's own. Assets: no retained span contains a figure environment, a graphics-inclusion command, a PDF-page inclusion, a table or a TikZ picture; no image file is copied into this package, the package declares no asset dependency and no image file is redistributed with it. Licence: version arXiv:2401.04220v2 of this article is distributed under the Creative Commons Attribution 4.0 International licence, as recorded in the arXiv licence metadata for this exact version and shown on the abstract page https://arxiv.org/abs/2401.04220v2. A full rights notice is delivered at licenses/arxiv-2401.04220-CC-BY-4.0.txt. Characters: every retained excerpt is pure ASCII, so the delivered engine, XeLaTeX with its default Latin Modern text font, reports no missing character.
\begin{proposition}\label{prop:reconstructshelling}
    Let $\cO$ be a good acyclic orientation of $G$. Assume $|V(G)| = n$. Let $t_1$ be the unique sink of $G$ with respect to $\cO$. Now define $t_2,\dots, t_n$ recursively by taking $t_i$ to be a sink of $G[V(G) \setminus \{t_1, \dots, t_{i-1}\})]$ for each $2 \le i \le n$. Then $T_1, \dots, T_n$ is a shelling of $\Delta$.
\end{proposition}

\begin{claim}\label{clm:initial}
For every $1 \le i \le n$, 
\begin{equation}\label{eqn:initialset}
   G[\cS_{k_i}(\rho_i)] \cup \cdots \cup G[\cS_{k_n}(\rho_n)] = G[\cV^\Delta_{k_i}(\rho_i)] \cup \cdots \cup G[\cV^\Delta_{k_n}(\rho_n)].
\end{equation}
\end{claim}

\begin{proof}[Proof of Claim \ref{clm:initial}]
Let $1 \le i \le n$. Let $G_{\cS, i}$ and $G_{\cV^\Delta, i}$ denote the graphs on either side of (\ref{eqn:initialset}) respectively. We first show that 
\[
G_{\cV^\Delta, i} = G[V(G) \setminus \{t_1, \dots, t_{i-1}\}].
\]
By the proof of Proposition \ref{prop:reconstructshelling}, if $t \in \cV^\Delta_{k_j}(\rho_j)$ for any $i \le j \le n$, then $T$ appears after $T_j$ in the shelling. Therefore, the two graphs have the same vertex set. Moreover, if $t_\ell \xrightarrow{\cO} t_j$ is a directed edge in $G$ with $\ell > j \ge i$, then $t_\ell \in \cV^\Delta_{k_j}(\rho_j)$, so the edge is also contained in $G[\cV^\Delta_{k_j}(\rho_j)]$. 

Let $\Phi_{k_i}(t_i)$ denote the principal $k_i$-frame representing $\rho_i$. It follows from the above that 
\[
G_{\cV^\Delta, i} = \Phi_{k_i}(t_i) \cup G_{\cV^\Delta, i+1}. 
\]
We now prove that $G_{\cS, i} = G_{\cV^\Delta, i}$ for every $1 \le i \le n$; we do so by backward induction on $i$. The base case is to show that $G[\cS_{k_n}(\rho_n)] = G[\cV^\Delta_{k_n}(\rho_n)]$. The restriction face of the last facet in any shelling of a sphere is the facet itself, so $k_n = 0$ and $G[\cS_{k_n}(\rho_n)] = G[\cV^\Delta_{k_n}(\rho_n)] = \{t_n\}$. Let $1 \le i \le n-1$ and assume inductively that $G_{\cS, i+1} = G_{\cV^\Delta, i+1}$. Then 
\[
G_{\cS, i} = G[\cS_{k_i}(\rho_i)] \cup G_{\cV^\Delta, i+1} \supseteq \Phi_{k_i}(t_i) \cup G_{\cV^\Delta, i+1} = G_{\cV^\Delta, i}. 
\]
On the other hand, observe that $V(G) \setminus \{t_1, \dots, t_{i-1}\}$ is initial in $G$. Since $\cS_{k_i}$ is star-like, the sink $t_i$ of $G[\cS_{k_i}(\rho_i)]$ is unique. Thus there is a directed path from every vertex in $\cS_{k_i}(\rho_i)$ to $t_i$. This implies $G[\cS_{k_i}(\rho_i)] \subseteq G[V(G) \setminus \{t_1, \dots, t_{i-1}\}]$. Thus
\[
G_{\cS, i} = G[\cS_{k_i}(\rho_i)] \cup G_{\cV^\Delta, i+1} \subseteq G_{\cV^\Delta, i}. 
\]
This completes the induction and (\ref{eqn:initialset}) follows. 
\end{proof}

\end{document}
